\documentclass[11pt]{article}
\usepackage[margin=1in]{geometry}
\usepackage[utf8]{inputenc}
\usepackage{amsmath,amssymb}
\usepackage{graphicx}
\usepackage{booktabs}
\usepackage{caption}
\usepackage[hidelinks]{hyperref}
\usepackage{enumitem}

\title{\bfseries What does $\Delta$KL measure?\\
\large Construct validity of the trademark resonance signal:\\
within-firm patent co-movement and external-list discriminant validity}
\author{Eric Silver}
\date{May 2026 --- working note}

\newcommand{\dkl}{$\Delta$KL}

\begin{document}
\maketitle

\begin{abstract}
\noindent For each USPTO trademark filing we score
$\Delta\mathrm{KL}=\mathrm{KL}(\text{filing}\,\|\,\text{past})-\mathrm{KL}(\text{filing}\,\|\,\text{future})$
on its goods/services text against same-class reference windows --- DeDeo
\emph{resonance}: positive when a filing's vocabulary is novel against the recent
past yet close to the near future (``the field moved toward it''). Does this
measure \emph{innovation}? Using the public-firm panels, three results sharpen
the construct. (1)~\dkl{} is essentially uncorrelated with patents in the cross
section (Pearson $+0.010$, Spearman $-0.015$) and, within a firm over time, a
\emph{mild substitute} for patenting ($-0.05\sigma$ per log-patent, $t\!\approx\!-8$,
strongest at the one-year patent lag) --- robust to filing-count and patent-matching
choices. (2)~Yet firms that BCG (2021--23) and MIT Technology Review (2015)
independently rank ``most innovative'' score $+0.29$ to $+0.50$ population-$\sigma$
\emph{higher} on \dkl{} ($p<0.012$; pooled $+0.34\sigma$, $p=0.0003$). (3)~The
headline ``$+28$pp per $\sigma$'' four-year debut-return effect is a small-sample
right-tail artifact: its horizon profile is non-monotone, peaking at 4y and
collapsing at 5y, and it runs $2.4\times\!\to\!9.3\times$ the firm-year estimate as
the horizon lengthens. Together these say \dkl{} is \emph{not} a proxy for
patent-style technical invention but \emph{does} track market/business-model
innovation as recognized by expert observers. The right question is not
``innovation versus churn'' but \emph{which} innovation construct \dkl{} names.
\end{abstract}

\section{The construct question}
The program's theory (PROPOSAL.md) defines innovation as \emph{novelty that
diffuses}: a business concept is ``innovative'' if its elements are
\emph{replicated} after the firm introduces them. The resonance statistic
operationalizes exactly this --- high \dkl{} marks language that was unusual at
filing and then adopted. On that theory, a filing being ahead of a class whose
vocabulary subsequently moves toward it is not a churn \emph{confound} to be
purged; it \emph{is} the innovation signal. The sharper question is therefore not
whether to subtract ``category churn'' out of \dkl{}, but whether \dkl{} names a
recognizable innovation \emph{construct} and how it relates to the established
ones. We test it against the two best external yardsticks available in the
repository: patents (technical invention) and expert ``most-innovative'' lists
(market/business-model innovation).

\paragraph{Scope.} The raw 10M-filing corpus is not on this branch, so the
corpus-level tests (term provenance, half-life sensitivity, ID-Manual
codification, attorney fixed effects) are deferred (\S\ref{sec:blocked}). All
results below run on the published firm-year panels and the external-list
crosswalk.

\section{Within-firm: \dkl{} and patenting are temporal substitutes}
Panel: 174{,}569 firm-years for 14{,}463 trademark owners name-matched to a
PatentsView assignee, collapsed to one row per (firm, year). We regress
standardized \dkl{} on $\log(1+\text{patents})$ with firm and year fixed effects
and firm-clustered standard errors; $\sigma(\Delta\mathrm{KL})=0.214$.

\begin{table}[h]\centering\small
\caption{Within-firm \dkl{} vs.\ patenting. Coefficients in \dkl{} standard
deviations per unit $\log(1+\text{patents})$.}
\begin{tabular}{lrrr}
\toprule
Specification & coef ($\sigma$) & SE & $t$ \\
\midrule
Between-firm correlation (firm means) & $-0.020$ & --- & --- \\
Within-firm correlation (firm+year demeaned) & $-0.015$ & --- & --- \\
Two-way FE, contemporaneous & $-0.048$ & $0.006$ & $-7.9$ \\
\quad robustness: patents = max over matches & $-0.048$ & $0.006$ & $-8.0$ \\
\quad controlling $\log(\text{n\_filings})$ & $-0.051$ & $0.006$ & $-8.4$ \\
Lead/lag: patents$_{t-1}$ & $-0.054$ & $0.009$ & $-5.8$ \\
\quad\quad\quad\ \ patents$_{t}$ & $-0.019$ & $0.010$ & $-1.9$ \\
\quad\quad\quad\ \ patents$_{t+1}$ & $+0.010$ & $0.010$ & $+1.0$ \\
\bottomrule
\end{tabular}
\end{table}

The coefficient is negative and precisely estimated, and \emph{strengthens} when
we control for the number of filings (which enters positively, $+0.055$,
$t=10.5$) --- ruling out the mechanical worry that high-patent years simply have
more filings whose averaging shrinks the mean toward the class. The effect is
largest at the one-year patent \emph{lag}: a burst of patenting is followed by
\emph{lower} trademark resonance the next year. Measurement noise from the fuzzy
owner$\to$assignee join attenuates toward zero, so the true within-firm relation
is, if anything, more negative.

\begin{figure}[h]\centering
\includegraphics[width=0.62\textwidth]{results/wsC_leadlag.png}
\caption{Within-firm \dkl{} declines with recent patenting; no positive
co-movement at any lead or lag.}
\end{figure}

\paragraph{Reading.} This decisively refutes \emph{convergent} validity with
patents --- not merely the near-zero cross-sectional correlation, but a genuine
within-firm \emph{opposition}. In the years a firm leans on patents it writes
more class-typical trademark language; in the years it pioneers commercial
vocabulary it patents less. Patents and \dkl{} capture different, partly
substituting modes of novelty.

\section{Discriminant validity against expert innovation lists}
If \dkl{} names an innovation construct at all, firms an independent panel calls
``most innovative'' should sit above the population. We compare the firm-mean
\dkl{} of list members (matched to a CIK in our public panel) against the
non-listed baseline (population: 7{,}431 firms, mean $-0.018$, $\sigma=0.175$).

\begin{table}[h]\centering\small
\caption{\dkl{} of expert-rated innovators vs.\ the non-listed baseline.}
\begin{tabular}{lrrrr}
\toprule
List & $n$ matched & listed mean & diff (pop-$\sigma$) & Mann--Whitney $p$ \\
\midrule
BCG Most Innovative (2021--23) & 31 & $+0.032$ & $+0.285$ & $0.012$ \\
MIT TR50 Smartest (2015)       & 17 & $+0.069$ & $+0.497$ & $0.0006$ \\
Pooled (union)                 & 43 & $+0.042$ & $+0.344$ & $0.0003$ \\
\midrule
\multicolumn{5}{l}{\textit{Convergent reference --- patents (174{,}569 firm-years):}}\\
Pearson (raw / log) & & $+0.010 / -0.009$ & & \\
Spearman (rank)     & & $-0.015$ & & $8\!\times\!10^{-10}$ \\
\bottomrule
\end{tabular}
\end{table}

\begin{figure}[h]\centering
\includegraphics[width=0.62\textwidth]{results/wsD_discriminant.png}
\caption{Population distribution of firm-mean \dkl{} with the BCG and MIT TR50
list means overlaid. Expert-rated innovators sit near the 67th--76th percentile.}
\end{figure}

\paragraph{Reading, with caveats.} \dkl{} aligns with expert, non-patent,
non-financial innovation judgments while being orthogonal-to-negative with
patents. That is the signature of a \emph{market/business-model} innovation
construct rather than a technical-invention one. Caveats are real and limit this
to suggestive: $n$ is small (only the US-public, CIK-matchable slice of global
lists matched), the comparison is each firm's all-years mean against a recent
list (no temporal alignment), and survivorship inflates the listed group. It is
nonetheless the first \emph{out-of-corpus} validation in the program --- the
lists were built without trademark text.

\section{The +28pp return effect is an artifact --- drop it}
The predictor is $z$-scored in \texttt{returns\_regression.py}, so the
$+0.284$ debut coefficient really is $+28$pp per $\sigma$. But it is not credible
as a headline: the debut excess-return profile is non-monotone and the debut /
firm-year ratio grows monotonically with horizon --- the fingerprint of a few
right-tail winners in a 1{,}447-firm sample, not a return premium.

\begin{table}[h]\centering\small
\caption{Excess-return coefficient on \dkl{} (per $\sigma$), by horizon.}
\begin{tabular}{lrrr}
\toprule
Horizon & Debut ($n\!\approx\!1.4$k) & Firm-year ($n\!\approx\!16$k) & ratio \\
\midrule
1y & $+0.036\ (t\,2.0)$ & $+0.015\ (t\,3.0)$ & $2.4\times$ \\
2y & $+0.067\ (t\,2.5)$ & $+0.035\ (t\,3.8)$ & $1.9\times$ \\
3y & $+0.192\ (t\,3.5)$ & $+0.050\ (t\,3.7)$ & $3.8\times$ \\
4y & $+0.284\ (t\,3.3)$ & $+0.041\ (t\,2.5)$ & $6.9\times$ \\
5y & $+0.149\ (t\,2.3)$ & $+0.016\ (t\,0.8)$ & $9.3\times$ \\
\bottomrule
\end{tabular}
\end{table}

\noindent Report the firm-year association ($+1.5$ to $+5$pp per $\sigma$, decaying
to insignificance by 5y) as the defensible financial fact, pending value-weighted,
delisting-adjusted, $5/95$-winsorized re-estimation of the debut panel.

\section{Discussion: which innovation construct?}
The ``innovation vs.\ churn'' dichotomy dissolves. \dkl{} is not a noisy patent
proxy contaminated by category churn; it is a measure of a \emph{different}
construct that happens to be near-orthogonal to patents. The evidence:

\begin{itemize}[nosep]
\item \dkl{} $\perp$ patents (cross-section) and $<0$ within-firm $\Rightarrow$
it does not track technical invention;
\item \dkl{} $> 0$ for BCG/MIT innovators $\Rightarrow$ it tracks market/
business-model innovation as expert observers perceive it.
\end{itemize}

\noindent On the ladder of names: the strong claim ``\dkl{} is \emph{the}
innovation measure'' is not licensed. The defensible primitive is \emph{lexical
resonance} --- a filing's forward-leaning position in its class vocabulary --- and
its validated economic content is \emph{market-facing innovation as a construct
distinct from invention}. The diffusion-as-innovation theory survives; what fails
is the assumption that this construct should converge with patents. It should
not, and does not.

\paragraph{Bridge to the diffusion program (sibling note).} A separate
Phase-0 test of \texttt{PROPOSAL\_diffusion.md} on the four-class
``Uber-for-X'' subset \{009, 042, 039, 035\} (2.32M granted filings,
\texttt{diffusion\_phase0\_note.tex}) now passes both §5 kill conditions:
the date/class shuffle null is rejected ($z=+46$ on the lag-1 autocorrelation
of LDA-theme trajectories), independently replicated on a separately-fitted
NMF panel ($z=+60.7$), and signed-ahead-precedes-behind is directionally
correct ($61\%$ of $(\theta, c)$ cells negative, modest magnitude consistent
with the $0.97$ collinearity between prospective and retrospective KL noted
above). The diffusion register --- themes as the unit, not the per-filing
scalar --- is the live continuation of the construct argument made here.

\paragraph{The localized innovation signal materializes (RQ6).} The diffusion
program's RQ6 --- does borrowed novelty pay at the firm level? --- is the
test this paper said would be the right place to look for an innovation
reading. With the SEC Financial Statement Data Sets crosswalked to USPTO
owners (7{,}233 exact matches, 4{,}999 CIK-linked firms in the borrowed-
novelty panel), the firm-year two-way FE regression of financial outcomes on
within-firm borrowed-gap finds: \textbf{contemporaneous gross margin
$+1.4$pp per $\sigma$ ($t=2.1$)} that fades by $t+1$, and \textbf{firms with
higher borrowed novelty have markedly lower R\&D intensity} cross-sectionally.
That mirror image --- borrowed novelty buys some margin in the year of the
filing while being a within-firm substitute for technical R\&D --- coheres
exactly with the within-firm $\Delta$KL/patent substitution result in §2:
borrowed novelty is the market-facing innovation channel; the firms who lean
on it lean less on the technical-invention channel. The localized innovation
reading is licensed on this margin-and-substitution shape, not on a global
convergent claim.

\section{Within-class diagnostics on a single technology corpus (Class 009)}

The firm-year results above use the aggregated panel. To stress-test the
construct's robustness at the filing level, we ran the full DeDeo flat-window
pipeline on a complete Class 009 (Software \& Electronics) corpus: 1{,}807{,}636
filings, 1893--2026, vocabulary $V = 151{,}193$ (min $df = 50$), per-filing
prospective + retrospective KL at $W \in \{3, 5, 7\}$ years. Four diagnostics
matter for what we can claim.

\paragraph{Pros/retr collinearity.} Across 1.56M scored filings,
$\mathrm{corr}(\mathrm{pros\_kl}, \mathrm{retr\_kl}) = \mathbf{+0.971}$. Signed
$\Delta$KL is a \emph{small residual of two large, highly correlated quantities}.
That is the point of the difference (it isolates timing from raw
distinctiveness) but also the fragility (small residuals are noisy at the
filing level).

\paragraph{Length artifact (F0b).} Mean $|\Delta\mathrm{KL}|$ falls
monotonically with filing length: $0.293$ for $<5$ terms, $0.265$ for 5--14,
$0.205$ for 15--39, $0.147$ for $\geq 40$. Short filings carry inflated
magnitudes by construction. Any per-filing analysis should down-weight or
control for length.

\paragraph{Window sensitivity (F0c).} $\Delta$KL \emph{magnitudes} are kernel-
sensitive, but the \emph{sign} is robust: only $10.4\%$ of filings flip sign
between $W=5$ and $W=3$ years and $6.9\%$ between $W=5$ and $W=7$, well under
the 20--30\% fragility threshold. The resonance direction (ahead vs.\ behind)
is not a kernel artifact.

\paragraph{Burn-in.} Mean signed $\Delta$KL at 1990 is $+2.05$ because the
pre-1990 software past is thin; from 1995 onward the per-year mean settles
within $\pm 0.07$ of zero. Restrict pre-1995 to a burn-in window.

\paragraph{Face validity post-burn-in.} Among 1996--2020 filings with $\geq 10$
terms, the top signed-$\Delta$KL marks are dominated by \textbf{2018
blockchain-era language}: STRAINCHAIN, CORPCOIN, MEDEXCHAIN, THE BLOCKCHAIN
BANKER, MEDXCHAIN, ART DEVICE. These are filings whose vocabulary was novel
at filing time and adopted afterward --- the construct firing on exactly what
it claims.

\paragraph{Term provenance (within-class Workstream D).} For each of 872{,}747
post-burn-in filings (1996--2020) we computed \texttt{novel\_share} = fraction
of its tokens whose first appearance in Class 009 was within $K=3$ years prior.
$\mathrm{corr}(\texttt{novel\_share}, \Delta\mathrm{KL}) = \mathbf{+0.29}$
(Pearson) and $+0.21$ (Spearman, $p \approx 0$). Filings with $\sim 0$ recently-
introduced class tokens average $\Delta\mathrm{KL} = -0.03$; filings with $\sim
5\%$ new tokens average $+0.10$. \textbf{The resonance signal is genuinely
new-vocabulary-driven within the class}, not recombination of established
terms (the corpus-new vs class-new split needs cross-class data; here we can
only confirm new-to-class drives the signal). The cross-class version on the
full diffusion subset is the natural extension and is covered in the sibling
note.

\section{What remains, and the decision rules}\label{sec:blocked}
These need the raw corpus (\texttt{data/processed/}) or external data not on this
branch. Each is scripted or specified; the decision rule is stated so the call is
made ex ante.
\begin{description}[nosep,style=nextline]
\item[Interaction estimand ($\beta\times$churn).] Needs per-NICE-class \dkl{}
and an \emph{independently constructed} churn measure (vocabulary turnover, not
the resonance kernel). Tests whether \dkl{}'s innovation content is moderated by
class dynamism. Rule: if the slope is positive only where churn $>0$, the
construct is intrinsically a firm$\times$class-phase quantity (consistent with the
diffusion theory), not a censored firm attribute.
\item[Term provenance (three-way).] Split \dkl{} into corpus-new, class-new
(recombination), and class-extant tokens. Rule: load on recombination $\Rightarrow$
rename the firm component ``recombinant novelty,'' not foresight.
\item[Half-life sensitivity.] Re-score sign of \dkl{} at $H\in\{1,2,4\}$y. Rule:
$>$20--30\% sign flips $\Rightarrow$ report \dkl{} as kernel-relative.
\item[ID-Manual codification (external).] Do high-\dkl{} filings use terms the
USPTO later standardizes, while those terms are still rare? This is the strongest
available external validation of the diffusion claim; needs the ID Manual version
history.
\item[Attorney fixed effects.] Needs an XML re-parse for correspondent of record.
Kills the ``drafting-fashion alone'' null; not a clean innovation extractor
(counsel sorting absorbs signal).
\item[External survival outcome.] Bankruptcy (PACER) or web-presence persistence
on a validation subsample, to break the trademark-maintenance circularity in the
survival U-shape.
\item[BISG demographics.] The ethnic-cluster work should replace surname-only
imputation with BISG using the owner addresses (verify the address is the owner's,
not the correspondent's, first).
\end{description}

\paragraph{Reproducibility.} \texttt{scripts/wsC\_within\_firm\_patents.py},
\texttt{wsD\_validity\_battery.py}, \texttt{wsE\_returns\_diagnostic.py},
\texttt{wsF\_make\_figs.py}; metrics in \texttt{paper/results/wsC\_*.json},
\texttt{wsD\_*.json}, \texttt{wsE\_*.json}. Run under \texttt{.venv-analysis}
(numpy/pandas/scipy/statsmodels/matplotlib).
\end{document}
